\chapter{Plan global}

\vfill
\localtableofcontents
\vfill
\newpage

\section{Introduction}

\begin{itemize}
    \item contexte : navigation autonome, pas de GNSS mais des champs physiques
    \item motivation : limites des méthodes classiques (Kalman, carte embarquée connue)
    \item objectif : reconstruction de la fonction d'observation par GP, intégration dans le cadre de la navigation bayésienne
    \item principales contributions
    \item annonce du plan
\end{itemize}

Ne pas oublier la progession des problèmes:
\begin{enumerate}
    \item drift of inertial navigation
    \item GNSS data blurred or unavailable
    \item high initial uncertainty
    \item non-injective, non-linear, multimodal environments
    \item undersampled, poorly-known environment model
    \item weight degeneracy, particle impoverishment
    \item non-stationarity of the field renders distant samples locally irrelevant
    \item high computational cost of matrix inversion and multiplications
    \item no preexisting data about the environment
\end{enumerate}

\section{Estimation bayésienne}

\begin{itemize}
    \item principes généraux et rappels de théorie de la mesure (cf cours de Nadia)
    \item filtre optimal
    \item filtres de Kalman et limites
    \item Kalman d'ensemble ?
\end{itemize}

\section{Filtrage particulaire}

\begin{itemize}
    \item idem pour les maths
    \item filtrage particulaire simple et limites
    \item filtres particulaire avancés : régularisé, auxiliaire, marginalisé
    \item filtre de Kalman sur RKHS (ailleurs ?)
\end{itemize}

\section{Processus gaussiens}

\begin{itemize}
    \item rappels de maths sur les processus stochastiques (s'inspirer Rasmussen)
    \item krigeage simple, ordinaire
    \item méthodes de réduction de rang
    \item lien avec les RKHS
\end{itemize}

\section{Krigeage sous-échantillonné}

\begin{itemize}
    \item erreurs théoriques
    \item heuristiques de sous-échantillonnage (fenêtrage, clusterisation,\\sous-échantillonnage aléatoire)
    \item parler de l'intérêt pour la navigation ? permet de reconstruire la fonction d'observation ?
    \item résultats numériques sans navigation
\end{itemize}

\section{CAK-PF}

\begin{itemize}
    \item fenêtrage, clustering par Mean-Shift
    \item apprentissage des hyperparamètres
    \item étude d'ablation
    \item précision en fonction du nombre de données disponibles
    \item résultats numériques (RMSE, temps de calcul, comparaisons avec la BCR)
\end{itemize}

\section{SLAM}

\begin{itemize}
    \item GP stationnaires ou stationnaires par morceaux
    \item multi capteurs en mailles impossibles car dimensions trop proches
    \item RBPF/méthode de Manon
    \item RKHS/méthode de Karim
\end{itemize}

\section{Ouverture sur la fusion de champs ?}

\begin{itemize}
    \item à voir mais information de Fisher théorique
    \item tests sur les 2 cas précédents
\end{itemize}

\section{Conclusion et perspectives}

\begin{itemize}
    \item résumé des contributions : GP + FP
    \item adaptation du modèle en fonction de l'état (carte dégradée)
    \item reconstruction à partir d'un dictionnaire de fonctions de base (sans carte)
    \item limites des méthodes
    \item perspectives futures (modèles, méthodologie, ...)
\end{itemize}
